\chapter{Before You Type}

Nova Fun \& Games is a book of small programs for NovaBASIC. It is meant to
feel like a program book you can work through at the keyboard: type a little,
run a little, change a little, and learn the machine by making it move.

This first volume sets up the shared framework and the first Nova-native
listings. Some programs are still short starter listings, but the shape is the
same: type a small program, run it, and then change it until it feels like
yours.

\section{The Disk}

The listings live on \diskname. The disk has no autoboot program. Mount it,
load a program, and run it from BASIC:

\begin{lstlisting}
LOAD "ROADHOG"
RUN
\end{lstlisting}

The image also includes \commonlisting, the shared code used by every program
in the book. The generated disk copies the common framework into each game file
so a reader can load and run a game directly. The separate \commonlisting\ file
is still present for the type-in workflow.

\begin{typeinbox}
If you are typing from the printed book, type \commonlisting\ once and save it.
Then type a game listing below it, starting at line 10. The game will call
line 60000 when it wants the common title screen.
\end{typeinbox}

\section{Variable Names}

NovaBASIC accepts longer variable names, but only the first two characters are
significant. In this book, program variables use one or two characters.

The shared framework owns names beginning with \texttt{Z}. Game listings should
avoid \texttt{Z*} variables unless they are setting the framework controls:
\texttt{ZT\$}, \texttt{ZA\$}, \texttt{ZS\$}, and \texttt{ZP}.

\section{Line Numbers}

The book uses a simple line-number map:

\begin{center}
\begin{tabularx}{0.88\textwidth}{>{\ttfamily}lX}
\toprule
10-9999 & Game code typed for the individual program.\\
50000-59999 & Reserved for later helper routines local to a large game.\\
60000-60280 & Shared Nova Fun \& Games common framework.\\
\bottomrule
\end{tabularx}
\end{center}

This keeps the program easy to list on screen and easy to repair after a typing
mistake.
